\section{The End of Patriarchy}

\begin{wrapfigure}{rt}{.3\textwidth}
\includegraphics[width=.3\textwidth]{Hypsipyle\_sauve\_Thoas\_BnF\_Francais\_599\_fol.\_16-300x296.jpg}
\end{wrapfigure}

\textbf{Giovanni Boccaccio}, in \emph{On Famous Women}, relates the story of the Queen of Lemnos. During the reign of her father, King Thoas, the women rebelled against the patriarchy. Boccaccio describes it:

\begin{quotex}
The women were seized by madness and withdrew their untamed necks from the yoke of men. Scorning the old king's rule, they unanimously decided that on the following night they would turn their knives against all the men. They did not fail to carry out their plan.
\end{quotex}


The women all carried out their plan, except for Hypsipyle. She secretly allowed her father to flee on a ship bound for Chios and pretended to cremate his corpse. She was then made queen. Boccaccio draws a moral from this story:

\begin{quotex}
The devotion of children toward their fathers is certainly very holy. What is more proper, more just, and more praiseworthy than to reward with humanness and honour those from whose labour we received nourishment when we were weak and who watched over us with solicitude, brought us to maturity with continuous love, taught us manners and gave us knowledge, enriched us with honours and abilities, and made us strong in morals and in intellect? Since Hypsipyle solicitously repaid this debt to her father, it is not without reason that she is placed among noble women.
\end{quotex}


Boccaccio is describing the fourth commandment to honour one's father and mother. As with all such commandments, a resistance develops. We don't know the reasons that motivated the women of Lemnos to murder their fathers, husbands, brothers, and sons. Yet we can be sure they felt justified, just as similar sentiments arise in our own time. Perhaps the fathers in Lemnos were not as ideal as those described by Boccaccio. That can't be the whole story, because a man who appears today to be such a father is still not protected against those knives.

I listened briefly tonight to an Internet figure who claims to be right wing, devoted to evidence and reason, before turning it off in contempt. He asserted that such devotion toward the fathers is an abnormality, based on genetic programming. Without any evidence or facts to back him up, he claimed that in early societies, the young had to respect and agree with their elders or risk losing the source of their sustenance. Those with the ``genes'' for conformity survived. This is a denial of a moral universe; as we see from the story, there is a moral choice involved. Piety is the virtue of loving one's parents and, by extension, one's ancestors and way of life; it is not a genetic defect.

To continue the story, \textbf{Jason} and his Argonauts arrived on the island. We can assume that Hypsipyle was filled with desire, so she invited Jason to her bed. The result was twin sons. She had to send them away under the laws of Lemnos, so she sent them to Chios to be raised by her father. This act revealed her earlier deception. The women revolted against her and she, too, had to escape on a ship. Things did not go well for Hypsipyle at this point. Pirates seized and enslaved her. She was then given to Lycurgus and served as a governess to his son. He was killed by a snake while in her care, bringing her great grief.

By chance, her sons found her and took her away to safety. Although she was cared for by her father and her sons, the other men she encountered were not so solicitous toward her. They used her for their own purposes.

Jason, in particular, who reneged on his promise of eternal fidelity to Hypsipyle, ended up in the eighth circle of hell for his deceit, seduction, and abandonment. (\textbf{Dante}, \emph{Inferno})


\flrightit{Posted on 2018-10-02 by Aeneas}

\begin{center}* * *\end{center}

\begin{footnotesize}
\begin{sffamily}
\texttt{argusandphoenix on 2018-10-04 at 07:36 said:}

Even assuming something comes ``after'' patriarchy, something improved, the era of the polar beings perhaps, it would necessarily follow that matriarchy is movement in the other direction. If we're stuck with patriarchy until we appreciate it, we will be in it a long, long time. The virulent section of the current Red Pill-manosphere crowd reminds me of Evola's warnings that an ``Assyrian'' or ``titanic'' period follows up the Demetrian Silver Age of Matriarchy (Revolt Against the Modern World).

\hfill

\end{sffamily}
\end{footnotesize}

